\hypertarget{report-on-use-of-ai}{%
\section{Report on Use of AI}\label{report-on-use-of-ai}}

During the writing of this paper, the team indeed utilized generative AI
large models for literature search, model introduction, and paper
structure optimization. The following details the team's use of AI, with
all dialogues in this report translated into English by the AI large
model.

\begin{longtable}[]{@{}l@{}}
\toprule
\endhead
\textbf{Q1：}Please help explain in detail what the Pareto Model is \\
\textbf{Output：}The Pareto Model, rooted in the Pareto Principle
(commonly known as the 80/20 Rule), is a theoretical and analytical
framework used to describe the skewed distribution of resources,
effects, or outcomes in various natural, social, and economic systems.
Named after Vilfredo Pareto, an Italian economist and sociologist who
first documented this phenomenon in the late 19th century, the model has
become a foundational tool for understanding imbalance and optimizing
decision-making. \#\#\# Core Concept The central tenet of the Pareto
Model is the **disproportionate relationship** between inputs and
outputs: a small fraction of factors (typically around 20\%, called the
"vital few") is responsible for a large majority of results (typically
around 80\%, called the "trivial many"). Importantly, the 80/20 ratio is
not a rigid mathematical constant but a heuristic---actual proportions
may vary (e.g., 70/30, 90/10) depending on the context, but the
underlying "skewness" remains consistent. For instance: - In wealth
distribution, 20\% of a population might hold 80\% of a country's total
wealth; - In business operations, 20\% of a company's products could
generate 80\% of its revenue; - In quality management, 20\% of defect
types may cause 80\% of production issues. \#\#\# Mathematical
Foundation: The Pareto Distribution The Pareto Model is mathematically
formalized through the **Pareto Distribution**, a continuous probability
distribution designed to characterize variables with a "heavy tail"
(i.e., a small number of extreme values and a large number of small
values). Its key components include: 1. **Probability Density Function
(PDF)**: \textbackslash{[} f(x) = \textbackslash begin\{cases\}
\textbackslash frac\{\textbackslash alpha
x\_0\^{}\textbackslash alpha\}\{x\^{}\{\textbackslash alpha + 1\}\} \&
\textbackslash text\{for \} x \textbackslash geq x\_0,
\textbackslash\textbackslash{} 0 \& \textbackslash text\{for \} x
\textless{} x\_0, \textbackslash end\{cases\} \textbackslash{]} -
\textbackslash( x \textbackslash): The random variable (e.g., income,
product sales, task impact); - \textbackslash( x\_0 \textbackslash): The
minimum value of \textbackslash( x \textbackslash) (scale parameter,
representing the threshold below which no observations exist); -
\textbackslash( \textbackslash alpha \textgreater{} 0 \textbackslash):
The shape parameter (Pareto index), which determines the degree of
skewness---smaller \textbackslash( \textbackslash alpha \textbackslash)
means a more uneven distribution. 2. **Cumulative Distribution Function
(CDF)**: \textbackslash{[} F(x) = 1 - \textbackslash left(
\textbackslash frac\{x\_0\}\{x\}
\textbackslash right)\^{}\textbackslash alpha \textbackslash quad
\textbackslash text\{for \} x \textbackslash geq x\_0, \textbackslash{]}
This function calculates the probability that \textbackslash( x
\textbackslash) is less than or equal to a given value, illustrating how
quickly the "vital few" dominate the distribution. \#\#\# Key
Characteristics 1. **Unevenness**: Rejects the assumption of linear
proportionality, emphasizing that outcomes are rarely distributed
equally across factors. 2. **Scalability**: Applies across scales (from
individual tasks to global systems) and domains (economics, business,
engineering, etc.). 3. **Prioritization Focus**: Highlights the value of
identifying and focusing on the "vital few" factors to maximize
efficiency---allocating resources to these factors yields
disproportionate returns. \#\#\# Practical Applications The Pareto
Model's simplicity and versatility make it widely used in diverse
fields: - **Economics \& Finance**: Analyzing income inequality, market
concentration, and wealth distribution (e.g., measuring the Gini
coefficient alongside Pareto metrics). - **Business Management**:
Optimizing product portfolios (prioritizing high-revenue items),
improving customer retention (focusing on high-value clients), and
reducing operational waste (targeting key inefficiencies). - **Quality
Control**: In lean manufacturing and Six Sigma, the "Pareto Chart" (a
hybrid of bar chart and line graph) identifies the most frequent causes
of defects, guiding process improvements. - **Productivity \& Project
Management**: Prioritizing tasks to focus on high-impact activities,
reducing time spent on low-value work, and allocating team resources
effectively. \#\#\# Limitations - **Descriptive, Not Prescriptive**: The
model describes *what is* (a skewed distribution) but not *why* it
occurs or how to change it---it must be paired with causal analysis for
actionable insights. - **Not Universal**: Not all phenomena follow a
Pareto distribution; some systems (e.g., normal distributions in natural
sciences) exhibit more balanced outcomes. - **Risk of
Oversimplification**: Overreliance on the 80/20 ratio may lead to
neglecting the "trivial many," which could become critical over time
(e.g., niche products that grow into major revenue streams). In essence,
the Pareto Model is a powerful tool for recognizing imbalance, enabling
individuals and organizations to prioritize effectively, optimize
resource allocation, and achieve more with less by focusing on the
factors that drive the most significant results. \\
\bottomrule
\end{longtable}

\begin{longtable}[]{@{}
  >{\raggedright\arraybackslash}p{(\columnwidth - 0\tabcolsep) * \real{1.00}}@{}}
\toprule
\endhead
\textbf{Q2：}What authoritative literatures are available on the
technical principles, feasibility of space elevators, and their
application in space colonization? \\
\textbf{Output：}The following literatures cover core aspects of space
elevators (e.g., technical design, economic paradigm, deployment
methods) and their integration with space colonization (e.g., Mars, Moon
colonization), with sources including the International Space Elevator
Consortium (ISEC), NASA, arXiv, and authoritative publishing platforms,
ensuring academic rigor and practical reference value:

1. Literatures Focused on Space Elevator Fundamentals \& Technical
Feasibility

Title: Space Elevator: A Revolutionary Earth-to-Space Transportation
SystemAuthors: Bradley Edwards, Eric WestlingPublisher/Year: BC Edwards
Publishing, 2002Core Content: This book is a foundational work on space
elevator technology. It systematically elaborates on the design
framework of space elevators (including tether materials, counterweight
configuration, and climber systems), and quantifies the efficiency
advantages of space elevators over traditional rockets (e.g., energy
consumption is only 1\% of that of rockets). It also proposes a phased
construction plan suitable for early technical conditions.Value:
Provides a comprehensive theoretical basis for understanding the basic
structure and technical logic of space elevators, and is widely cited in
subsequent space elevator research.

Title: "Changing the Economic Paradigm for Building a Space
Elevator"Authors: Kevin Barry, Eduardo Pineda AlfaroConference/Year:
72nd International Astronautical Congress (IAC-21), Dubai, UAE, October
2021Core Content: The paper breaks through the traditional "high-cost
construction" cognition of space elevators, analyzes the economic
feasibility of space elevators from the perspective of industrial chain
collaboration (e.g., shared use of ground bases, batch production of
tether materials), and predicts that the cost of space elevator
transportation can be reduced to \$100-200 per kilogram in the mature
stage.Value: Supplements the economic dimension of space elevator
research, which is crucial for evaluating the cost-effectiveness of
space colonization transportation.

Title: "Non-Equatorial Uniform-Stress Space Elevators"Author: Blaise
GassendPlatform/Year: arXiv (2509.16231), October 2025Core Content: This
paper expands the design scope of space elevators beyond the traditional
equatorial anchor. It demonstrates the feasibility of non-equatorial
space elevators through mechanical analysis (equilibrium of gravity,
centrifugal force, and tether tension), and points out their advantages:
avoiding crowded geosynchronous orbits (GEO) and high-radiation Van
Allen belts, which is more suitable for human-carrying transportation in
space colonization.Value: Provides a new design direction for space
elevator anchor site selection, especially for space colonies that need
to avoid specific orbital resources.

2. Literatures on Space Elevators \& Space Colonization Application

Title: Space Colonization Using Space-Elevators from Phobos
(PDF)Institution/Year: NASA (NTRS Citation: 20030065879), 2003Core
Content: This is a classic NASA study on Mars colonization. It proposes
a "Mars-Phobos space elevator system": using Phobos (a moon of Mars) as
a raw material source and tether anchor, deploying two space
elevators---one connecting Phobos to the upper atmosphere of Mars (for
receiving surface payloads) and the other extending outward from Phobos
(for transporting materials to the Earth-Moon system or asteroid belt).
It also analyzes the role of this system in supporting Mars colony
infrastructure construction and space industry development.Value:
Provides a typical case of integrating space elevators with planetary
colonization, which can be referenced for designing Moon colony
transportation systems (e.g., using lunar satellites as elevator
anchors).

Title: "Exponential Tethers for Accelerated Space Elevator
Deployment"Author: Blaise GassendPlatform/Year: arXiv (2412.17198),
December 2024Core Content: The paper proposes an "exponential tether"
design for space elevators---by making the tether cross-section vary
exponentially with altitude, the tether can be reeled in and out at both
ends without changing its stress profile. This design significantly
shortens the deployment cycle of space elevators (reducing the time from
initial launch to full operation by 40\% compared to traditional
methods) and is suitable for rapid transportation support in the early
stage of space colony construction.Value: Solves the problem of "slow
deployment of space elevators" in space colonization, and provides a
technical path for timely material supply to colonies.

3. Literatures on Space Elevator History \& Industry Practice

Title: Space Elevators: A HistoryAuthor: David RaittReport/Year: ISEC
Report, 2017Core Content: Compiled by the International Space Elevator
Consortium (ISEC), this report sorts out the development history of
space elevators from the 19th century (Tsiolkovsky's "sky ladder"
concept) to the 21st century (carbon nanotube material breakthroughs).
It also summarizes global research projects (e.g., Japan's Obayashi
Corporation's space elevator plan, NASA's advanced technology research)
and industry cooperation models.Value: Helps understand the evolution of
space elevator technology and the current status of industrial research,
and provides a reference for the "maturity of space elevator technology"
in space colonization planning.

Title: Road to the Space Elevator EraAuthors: Peter Swan, David Raitt,
John Knapman, et al.Publisher/Year: Virginia Edition Publishing Company
(Science Deck), 2019 (ISBN: 978-0-9913370-3-3)Core Content: This book
integrates technical research and practical planning, covering topics
such as space elevator material upgrades (from carbon nanotubes to
graphene composites), space debris avoidance strategies, and the
integration of space elevators with space habitats (e.g., connecting to
rotating artificial gravity habitats). It also predicts that space
elevators will become the core infrastructure for lunar and Martian
colonization by the 2050s.Value: Bridges the gap between space elevator
technology and space colonization practice, and is highly relevant to
the 2050 Moon colony construction scenario in the 2026 MCM Problem B.

4. Supplementary Literature on Space Elevator Materials \& Feasibility

Title: "太空升降舱\_百科" (Space Elevator Encyclopedia)Platform/Year:
抖音百科 (Douyin Encyclopedia), February 2025Core Content: This
Chinese-language reference systematically introduces the composition of
space elevators (counterweight, tether, ground base, climber), material
requirements (carbon nanotube's strength-to-weight ratio advantage), and
global research progress (e.g., Japan's estimated \$93 billion
construction cost, NASA's tether durability tests). It also analyzes
technical challenges (batch production of carbon nanotubes, space debris
impact prevention).Value: Provides a concise and easy-to-understand
overview of space elevator basics, suitable for quickly grasping key
technical indicators and feasibility constraints. \\
\bottomrule
\end{longtable}

\begin{longtable}[]{@{}
  >{\raggedright\arraybackslash}p{(\columnwidth - 0\tabcolsep) * \real{1.00}}@{}}
\toprule
\endhead
\textbf{Q3：}What are the current levels of launch costs for mainstream
rockets worldwide, and what authoritative data supports them? \\
A:The current global rocket launch costs show significant differences
due to factors such as reusability, rocket type, and regional technology
gaps. Authoritative data mainly comes from NASA reports, industry
analysis institutions (e.g., TrendForce), and enterprise cost
disclosures, with specific details as follows:

International mainstream rockets (taking SpaceX as a benchmark):

According to NASA's 2018 report (Summary 1) and 2020 report (Summary 5),
NASA's Space Shuttle had a launch cost of approximately \$54,500 per
kilogram for delivering 27,500 kg to Low Earth Orbit (LEO). In contrast,
SpaceX's Falcon 9 now advertises a cost of \$2,720 per kilogram for
launching 22,800 kg to LEO, a 20-fold reduction in cost.

For heavy-lift tasks, Summary 6 notes that SpaceX's Falcon Heavy has a
unit cost of \$1,400--1,500 per kilogram for LEO missions, benefiting
from scale effects of larger payload capacity.

Cost breakdown data (Summary 3) shows that the total cost of a new
Falcon 9 is \$50 million; when reused 10 times, the single-launch cost
drops to \$18 million, with the unit cost further optimized to around
\$2,700 per kilogram.

Domestic Chinese rockets (divided into "national team" and private
enterprises):

"National team" rockets (e.g., Long March series): As disposable
rockets, their costs are relatively rigid. Summary 6 states that the
Long March 2D has a launch cost of 50,000--80,000 RMB per kilogram (some
sources reach nearly 100,000 RMB/kg), the Long March 6A costs about
50,000 RMB per kilogram, and the Long March 5B (for large LEO payloads)
costs approximately 80,000+ RMB per kilogram.

Private commercial rockets: Summary 2 and 6 indicate that their current
market quotations range from 50,000--100,000 RMB per kilogram. In bids
for projects like the "Qianfan Constellation," enterprises such as
Tianbing Technology and LandSpace have offered winning bids of around
30,000 RMB per kilogram, with aggressive target prices of 20,000 RMB per
kilogram (based on future reusability expectations), though actual costs
remain around 50,000 RMB per kilogram due to immature technology.

Q: What is the mechanism by which rocket recovery technology reduces
launch costs, and what are the actual cost-reduction cases?A: Rocket
recovery technology reduces costs by "reusing high-value components" and
"spreading fixed costs across multiple launches," with verified effects
supported by both international and domestic cases:

Core cost-reduction mechanism:Summary 2 and 3 point out that the
first-stage rocket (including engines and airframe structure) accounts
for 60\%--70\% of a rocket's total cost. Recovering and refurbishing
these core components transforms "one-time hardware investment" into
"multi-time shared costs," directly cutting the largest single cost
item. Additionally, reusable rockets reduce the need for manufacturing
new rockets, lowering fixed costs such as mold opening and material
procurement.

Actual cases:

International (SpaceX Falcon 9): Summary 2 notes that after achieving
first-stage recovery and reuse, the single-launch cost of Falcon 9
dropped from approximately \$50 million (new rocket) to \$15--30
million; when a booster is reused 10 times, the single cost falls to
about \$18 million (a 64\% reduction from the new rocket cost). Summary
3 further specifies that with 15 reuses, the single-launch cost is
compressed to \$17 million, with a gross profit margin of 74\%.

Domestic (Chinese enterprises): Summary 2 states that LandSpace plans to
provide "used rocket" commercial launches in Q4 2026, targeting a 40\%
cost reduction; Zhongke Aerospace expects its Lijian-2 recoverable
rocket to lower the cost from 60,000--70,000 RMB per kilogram to over
20,000 RMB per kilogram.

Q: What gaps exist in launch cost control between domestic and
international commercial aerospace sectors, and what are the catch-up
goals of domestic enterprises?A: There is a "disruptive gap" in launch
cost control between domestic and international commercial aerospace,
with domestic enterprises setting clear phased catch-up goals:

Current gaps:According to Summary 6, the international leading level
(represented by SpaceX) has a unit cost of less than \$3,000 per
kilogram (≈21,000 RMB/kg) under reuse mode, while domestic launch costs
are generally 50,000--100,000 RMB per kilogram, approximately 2--3 times
higher. Specifically:

International: SpaceX's Falcon 9 has an actual marginal cost of
\textasciitilde\$2,720 per kilogram (Summary 1, 5), and Falcon Heavy's
unit cost is as low as \$1,400--1,500 per kilogram (Summary 6).

Domestic: Disposable "national team" rockets (e.g., Long March 2D) cost
50,000--100,000 RMB per kilogram; private enterprise quotations are
50,000--100,000 RMB per kilogram, with "futuristic prices" of
20,000--30,000 RMB per kilogram based on reusability expectations (not
yet realized in practice, Summary 6).

Domestic catch-up goals:

Short-term (2026--2027): Summary 2 and 6 note that LandSpace, Tianbing
Technology, and other enterprises aim to reduce launch costs to below
30,000 RMB per kilogram by 2027 through mature recoverable rockets,
achieving a nearly 50\% cost reduction.

Medium-term (2030): With the large-scale deployment of constellations
(e.g., China's "Guowang" Constellation), domestic commercial rockets are
expected to reach the 20,000--30,000 RMB per kilogram range, gradually
narrowing the gap with SpaceX's current level (Summary 6).

Q: What is the future trend of rocket launch cost reduction, and what
technologies or models can drive further cost cuts?A: The future of
rocket launch costs will show a "phased downward" trend, driven by three
core forces: reusable technology iteration, advanced propulsion
technology, and business model innovation:

Short-to-medium term (2026--2030): Cost reduction centered on "partial
reuse"As summarized in 2, 3, and 6, domestic recoverable rockets (e.g.,
LandSpace Zhuque-3, Tianbing Technology Tianlong-3) will achieve stable
reuse of first-stage rockets, reducing costs to 20,000--30,000 RMB per
kilogram domestically and maintaining SpaceX's international leading
cost at \textasciitilde\$2,000 per kilogram. Additionally, "rideshare"
models (Summary 3) will further分摊 costs---for example, SpaceX's
Transporter missions reduce the per-kilogram cost to below \$2,720 by
sharing launch capacity among multiple customers.

Long-term (after 2030): Cost revolution driven by "full reuse" and
advanced propulsion

Full-reusable rockets: Summary 4 and 6 mention that SpaceX's Starship (a
fully reusable super-heavy launch vehicle) targets a launch cost of
\$100 per kilogram (or even \$67 per kilogram in ideal operations),
which would 颠覆性 ly lower costs. Domestic enterprises are also
exploring full-reuse technologies, with a long-term goal of "affordable
space transportation" (Summary 6).

Advanced propulsion technologies: Summary 4 notes that technologies such
as electric propulsion (high specific impulse, low fuel consumption),
laser propulsion, and solar sail propulsion will reduce reliance on
chemical propellants (which account for over 90\% of a rocket's mass)
and further cut fuel and launch mass costs. For example, electric
propulsion has been used in over 200 satellites for orbit maintenance,
and its application in launch vehicles is expected to reduce fuel costs
by 30\% (Summary 4).

Q: How does the reduction in launch costs affect the planning of space
missions (e.g., lunar bases, space stations)?A: The reduction in launch
costs has a transformative impact on space mission planning, mainly
reflected in optimizing supply methods, improving mission feasibility,
and reducing life-cycle costs:

Optimizing life support supply methods for long-duration
missions:Summary 1 (NASA report) points out that high launch costs
previously forced long-duration missions (e.g., the International Space
Station, ISS) to rely on water and oxygen recycling systems to reduce
logistics mass. However, with the new low launch costs (e.g., Falcon 9's
\$2,720 per kilogram), "open-loop life support" (direct resupply of
water and oxygen) has become cheaper than recycling. For future missions
like long-term lunar bases, direct logistics resupply will be more
cost-effective than building complex recycling systems (Summary 1).

Reducing supply costs for large-scale space infrastructure:Summary 5
states that the cost of resupplying the ISS has been reduced by a factor
of 4 due to commercial launch technologies (e.g., Falcon 9 + Dragon
capsule). For the 100,000-person lunar colony in the 2026 MCM Problem B,
lower launch costs will directly cut the cost of transporting 100
million metric tons of materials and subsequent water supplies---for
example, if the unit cost drops from \$12,000 per ton (traditional
rockets) to \$3,000 per ton (reusable rockets), the total material
transportation cost will decrease by 75\%.

Enabling large-scale space missions:Summary 2 and 6 note that when
launch costs fall to \$3,000 per kilogram or below, previously
unaffordable missions (e.g., large LEO constellations with 10,000+
satellites, long-term Mars bases) become commercially viable. For lunar
colonies, lower costs also support more frequent crew rotations and the
transportation of non-essential materials (e.g., tourist facilities),
promoting the transition from "survival-oriented" to "livable-oriented"
colonies. \\
\bottomrule
\end{longtable}

\begin{longtable}[]{@{}
  >{\raggedright\arraybackslash}p{(\columnwidth - 0\tabcolsep) * \real{1.00}}@{}}
\toprule
\endhead
\textbf{Q4：}Could you demonstrate a complete, runnable code example for
solving practical multi-objective optimization problems using the Pareto
Model, with detailed explanations of key logic and steps? \\
\textbf{A:} Below is a practical code example that uses the Pareto Model
to solve a **two-objective production scheduling optimization problem**
(minimizing production cost + minimizing production time). The code is
based on Python (no complex third-party optimization libraries, easy to
understand), and includes core steps such as candidate solution
generation, Pareto optimal solution screening, and visualization of the
Pareto frontier.

\#\#\# Core Scenario Definition

A factory produces two types of products (Product A and Product B) with
the following constraints and objectives:

- **Constraints**:

1. Total raw material usage ≤ 100 units (Product A consumes 2
units/piece, Product B consumes 3 units/piece);

2. Total production capacity ≤ 80 hours (Product A takes 1 hour/piece,
Product B takes 2 hours/piece);

3. Production quantity of each product ≥ 0 (integer).

- **Objectives**:

1. Minimize total cost: Cost = 5×A + 8×B (Product A costs \$5/piece,
Product B costs \$8/piece);

2. Minimize total production time: Time = 1×A + 2×B (consistent with
capacity constraints, but treated as an independent objective).

\#\#\# Complete Code Implementation

```python

import numpy as np

import pandas as pd

import matplotlib.pyplot as plt

\# ==========================================

\# Step 1: Define problem parameters

\# ==========================================

\# Resource constraints

MAX\_MATERIAL = 100 \# Maximum raw material (units)

MAX\_CAPACITY = 80 \# Maximum production capacity (hours)

\# Cost and time coefficients for each product

COST\_COEFF = \{"A": 5, "B": 8\} \# Cost per unit (\$)

TIME\_COEFF = \{"A": 1, "B": 2\} \# Time per unit (hours)

\# ==========================================

\# Step 2: Generate all feasible candidate solutions

\# ==========================================

def generate\_feasible\_solutions():

feasible\_solutions = {[}{]}

\# Iterate over possible production quantities of Product A (integer)

for a in range(0, int(MAX\_CAPACITY / TIME\_COEFF{[}"A"{]}) + 1):

\# Calculate maximum possible B under material and capacity constraints

max\_b\_by\_material = (MAX\_MATERIAL - COST\_COEFF{[}"A"{]} * a) //
COST\_COEFF{[}"B"{]}

max\_b\_by\_capacity = (MAX\_CAPACITY - TIME\_COEFF{[}"A"{]} * a) //
TIME\_COEFF{[}"B"{]}

max\_b = min(max\_b\_by\_material, max\_b\_by\_capacity)

\# Iterate over feasible B quantities

for b in range(0, max\_b + 1):

\# Skip invalid solutions (negative B or exceeding constraints)

if (2*a + 3*b) \textgreater{} MAX\_MATERIAL or (1*a + 2*b)
\textgreater{} MAX\_CAPACITY:

continue

\# Calculate objectives (cost and time)

total\_cost = COST\_COEFF{[}"A"{]} * a + COST\_COEFF{[}"B"{]} * b

total\_time = TIME\_COEFF{[}"A"{]} * a + TIME\_COEFF{[}"B"{]} * b

feasible\_solutions.append(\{

"A\_quantity": a,

"B\_quantity": b,

"total\_cost": total\_cost,

"total\_time": total\_time

\})

\# Convert to DataFrame for easy processing

return pd.DataFrame(feasible\_solutions)

\# ==========================================

\# Step 3: Screen Pareto optimal solutions (core logic of Pareto Model)

\# Definition of Pareto optimality: A solution is Pareto optimal if
there is no other solution

\# that is better than it in all objectives, and at least one objective
is strictly better.

\# ==========================================

def get\_pareto\_optimal\_solutions(df):

pareto\_mask = np.ones(len(df), dtype=bool) \# Initialize all solutions
as Pareto candidates

for i in range(len(df)):

\# Skip if current solution is already marked as non-Pareto

if not pareto\_mask{[}i{]}:

continue

\# Current solution's objectives

current\_cost = df.iloc{[}i{]}{[}"total\_cost"{]}

current\_time = df.iloc{[}i{]}{[}"total\_time"{]}

\# Check if there exists another solution that dominates the current one

for j in range(len(df)):

if i == j:

continue \# Skip self-comparison

other\_cost = df.iloc{[}j{]}{[}"total\_cost"{]}

other\_time = df.iloc{[}j{]}{[}"total\_time"{]}

\# Dominance condition: other solution has lower/equal cost AND
lower/equal time, and at least one is strictly lower

if (other\_cost \textless= current\_cost) and (other\_time \textless=
current\_time) and (other\_cost \textless{} current\_cost or other\_time
\textless{} current\_time):

pareto\_mask{[}i{]} = False \# Mark current solution as non-Pareto

break \# No need to check other solutions

\# Filter Pareto optimal solutions

pareto\_solutions = df{[}pareto\_mask{]}.reset\_index(drop=True)

return pareto\_solutions

\# ==========================================

\# Step 4: Visualize Pareto frontier and results

\# ==========================================

def visualize\_pareto(df\_all, df\_pareto):

plt.figure(figsize=(10, 6))

\# Plot all feasible solutions (gray dots)

plt.scatter(df\_all{[}"total\_time"{]}, df\_all{[}"total\_cost"{]},

color="lightgray", alpha=0.6, label="Feasible Solutions")

\# Plot Pareto optimal solutions (red dots with black edges)

plt.scatter(df\_pareto{[}"total\_time"{]},
df\_pareto{[}"total\_cost"{]},

color="red", s=100, edgecolors="black", label="Pareto Optimal
Solutions")

\# Connect Pareto points to form the Pareto frontier

pareto\_sorted = df\_pareto.sort\_values(by="total\_time")

plt.plot(pareto\_sorted{[}"total\_time"{]},
pareto\_sorted{[}"total\_cost"{]},

color="darkred", linestyle="-\/-", linewidth=2, label="Pareto Frontier")

\# Add labels and title

plt.xlabel("Total Production Time (hours)", fontsize=12)

plt.ylabel("Total Production Cost (\$)", fontsize=12)

plt.title("Pareto Optimization: Production Cost vs. Production Time",
fontsize=14, fontweight="bold")

plt.grid(True, linestyle="-\/-", alpha=0.5)

plt.legend(fontsize=10)

plt.tight\_layout()

plt.show()

\# ==========================================

\# Step 5: Main execution logic

\# ==========================================

if \_\_name\_\_ == "\_\_main\_\_":

\# Generate all feasible solutions

all\_solutions = generate\_feasible\_solutions()

print(f"Total number of feasible solutions: \{len(all\_solutions)\}")

\# Get Pareto optimal solutions

pareto\_solutions = get\_pareto\_optimal\_solutions(all\_solutions)

print("\textbackslash nPareto Optimal Solutions (top 5):")

print(pareto\_solutions.head())

\# Visualize results

visualize\_pareto(all\_solutions, pareto\_solutions)

\# Output key Pareto solutions for decision-making

print("\textbackslash nKey Pareto Optimal Solutions for
Decision-Making:")

\# 1. Minimum cost solution

min\_cost\_sol =
pareto\_solutions.loc{[}pareto\_solutions{[}"total\_cost"{]}.idxmin(){]}

\# 2. Minimum time solution

min\_time\_sol =
pareto\_solutions.loc{[}pareto\_solutions{[}"total\_time"{]}.idxmin(){]}

\# 3. Balanced solution (closest to the ideal point: min cost + min
time)

ideal\_cost = pareto\_solutions{[}"total\_cost"{]}.min()

ideal\_time = pareto\_solutions{[}"total\_time"{]}.min()

pareto\_solutions{[}"distance\_to\_ideal"{]} = np.sqrt(

(pareto\_solutions{[}"total\_cost"{]} - ideal\_cost)**2 +

(pareto\_solutions{[}"total\_time"{]} - ideal\_time)**2

)

balanced\_sol =
pareto\_solutions.loc{[}pareto\_solutions{[}"distance\_to\_ideal"{]}.idxmin(){]}

print(f"\textbackslash n1. Minimum Cost Solution:")

print(f" Product A: \{min\_cost\_sol{[}'A\_quantity'{]}\} pieces,
Product B: \{min\_cost\_sol{[}'B\_quantity'{]}\} pieces")

print(f" Cost: \$\{min\_cost\_sol{[}'total\_cost'{]}\}, Time:
\{min\_cost\_sol{[}'total\_time'{]}\} hours")

print(f"\textbackslash n2. Minimum Time Solution:")

print(f" Product A: \{min\_time\_sol{[}'A\_quantity'{]}\} pieces,
Product B: \{min\_time\_sol{[}'B\_quantity'{]}\} pieces")

print(f" Cost: \$\{min\_time\_sol{[}'total\_cost'{]}\}, Time:
\{min\_time\_sol{[}'total\_time'{]}\} hours")

print(f"\textbackslash n3. Balanced Solution (closest to ideal point):")

print(f" Product A: \{balanced\_sol{[}'A\_quantity'{]}\} pieces, Product
B: \{balanced\_sol{[}'B\_quantity'{]}\} pieces")

print(f" Cost: \$\{balanced\_sol{[}'total\_cost'{]}\}, Time:
\{balanced\_sol{[}'total\_time'{]}\} hours")

```

\#\#\# Key Explanations of the Code

1. **Feasible Solution Generation**:

- Iterate over all integer quantities of Product A within capacity
limits, then calculate the maximum feasible quantity of Product B under
raw material and capacity constraints.

- This ensures all generated solutions comply with practical production
constraints (no invalid solutions).

2. **Pareto Optimal Solution Screening (Core Logic)**:

- Use a double-loop to compare each solution with all others: if
Solution B has **lower/equal cost AND lower/equal time** (and at least
one objective is strictly better), Solution A is marked as non-Pareto
optimal.

- The `pareto\_mask` array efficiently marks Pareto/non-Pareto solutions
to avoid redundant calculations.

3. **Visualization of Pareto Frontier**:

- All feasible solutions are plotted in gray, and Pareto optimal
solutions are highlighted in red---this直观地 shows the "trade-off
curve" between cost and time (the Pareto frontier).

- The frontier clearly illustrates that reducing cost inevitably
increases time (and vice versa), which is the core of multi-objective
optimization.

4. **Decision-Making Based on Pareto Solutions**:

- Extract three typical Pareto solutions for practical decision-making:

- Minimum cost solution (prioritize cost control);

- Minimum time solution (prioritize production efficiency);

- Balanced solution (closest to the "ideal point" of minimum cost +
minimum time, calculated via Euclidean distance).

\#\#\# Running Requirements \& Results

- **Dependencies**: Install required libraries with `pip install numpy
pandas matplotlib`.

- **Output Results**:

1. Total number of feasible solutions (typically \textasciitilde200+ for
this scenario);

2. A table of Pareto optimal solutions (10--15 solutions);

3. A scatter plot of the Pareto frontier;

4. Three key decision-making solutions with specific production
quantities, costs, and time.

This code is highly extensible---you can adapt it to the lunar colony
transportation optimization problem (replace "production quantity" with
"transport volume of space elevator/rocket", and "cost/time" with "total
transportation cost/total time") by modifying parameters and
constraints. \\
\bottomrule
\end{longtable}
